\section{The Contraction (Tzimtzum)}

What the plain story called the drawing-back, Kabbalah calls \textbf{tzimtzum}. It
is the founding move of the Lurianic system.

\subsection{Tzimtzum, the Contraction}

\textbf{Tzimtzum} means \enquote{contraction} or \enquote{constriction.} The problem
it solves is logical. If the infinite light, the Or Ein Sof, fills all reality with
no end, there is no room for any finite thing. A finite world cannot exist inside an
undifferentiated infinity, because the infinite would overwhelm it utterly.

The answer of Isaac Luria, the Ari of sixteenth-century Safed, is the tzimtzum. The
infinite light withdraws or contracts from a central point, leaving a vacated space.
Into that space a thin ray of light then re-enters, and the worlds are built within
it.

\begin{itemize}
 \item Without contraction there is no finitude. Infinity must make room.
 \item The contraction leaves a residue, so the space is not truly empty of God.
 \item A measured ray then enters the space to begin building worlds.
\end{itemize}

The tzimtzum is the hinge between the infinite and everything else. It is how the
unlimited becomes limited.

\subsection{The Empty Space, the Residue, the Ray}

Three terms name the parts of the contraction.

\textbf{The empty space.} The withdrawal opens a vacated region, the \textbf{chalal}
or makom panui. For the first time there is a \enquote{place} not saturated with
infinite light, a place where finite being can exist.

\textbf{The residue.} The light does not leave an absolute void. A faint trace stays
behind, the \textbf{reshimu}, the \enquote{impression.} It is the memory of the
Infinite within the finite, the seed on which the new emanation is built. The image
is a scent left in an emptied flask, or warmth in a room after the fire is out. The
light is gone, but something of it lingers and makes the next stage possible.

\textbf{The ray.} Into the vacated space the Infinite sends a single thin ray, the
\textbf{kav}, the \enquote{line} or \enquote{measure.} Unlike the unmeasured fullness
before the contraction, the kav is controlled and measured. All structure to come is
built along and around this line. The kav is the channel of emanation into the void.

\subsection{Literal or Figurative, the Great Debate}

How should the tzimtzum be understood? This is one of the most contested questions in
Jewish thought, and it deserves to be reported plainly, because the two readings
lead to two different pictures of reality.

\textbf{The literal reading}, tzimtzum kipshuto, holds that the contraction is real.
God genuinely withdrew the light from the vacated space. In that space the divine
presence is truly absent, or at least truly hidden in actuality. This is associated
with the Vilna Gaon and the Lithuanian Mitnagdic tradition.

\textbf{The figurative reading} holds that the contraction is not a real withdrawal.
God cannot literally absent himself, since he is everywhere and unchanging. The
tzimtzum is a concealment, not a departure. The light is hidden only from the
perspective of the worlds. From God's own view nothing has changed, and the Infinite
still fills the vacated space completely. This is the Hasidic position, developed most
sharply in Chabad.

\begin{longtable}{L{0.16\linewidth} L{0.36\linewidth} L{0.36\linewidth}}
\caption{The two readings of the contraction.}\\
\toprule
 & \textbf{Literal (kipshuto)} & \textbf{Figurative} \\
\midrule
The act & Real withdrawal & Concealment only \\
The space & Genuinely empty of light & God still fully fills it \\
Affects & The light itself & Only its revelation \\
Stream & Vilna Gaon, Mitnagdic & Hasidic, especially Chabad \\
\bottomrule
\end{longtable}

The labels are poles of a spectrum, not airtight boxes. The Vilna Gaon's own students
argued his view was more subtle than crude literalism, and the two camps are not as
far apart as polemics suggested. The fork is real all the same. If the tzimtzum is
literal, the world is a genuine domain where God's presence is withdrawn, so
independent existence and free will have firm ground. If it is figurative, the
world's separateness is an illusion of perspective, and only God truly exists. Almost
every other question, free will, prayer, evil, the value of the physical, follows
from this fork.

\subsection{The Purpose of Creation}

Why contract at all? Why is there a world?

Ein Sof has no need. It lacks nothing. So creation is not driven by deficiency. The
classical answer is the \textbf{divine will to bestow good} (ratzon). It is the nature
of the good to do good. The will to give is the root of creation.

This frames the whole cosmos as a \textbf{gift structure}. There is a giver who
wishes to bestow, and so there must be a receiver to bestow upon. The worlds exist so
that there can be a recipient of the divine good.

\subsection{The Bread of Shame}

There is a further constraint that shapes the whole design. A gift received for
nothing, with no effort, carries a hidden shame. The receiver feels the discomfort of
unearned charity, like eating bread one did not work for. This is the \textbf{bread
of shame} (nahama dekisufa). Pure receiving, with no giving back, is
degrading.

So if God's purpose is to bestow the greatest good, God cannot simply hand it over.
That would leave the receiver in the position of bread of shame. Instead God creates a
being that can \textbf{earn} its good through effort, choice, and labor. The earned
good is free of shame. It is the receiver's own.

\begin{itemize}
 \item God wishes to give the ultimate good.
 \item A freely given gift produces bread of shame in the receiver.
 \item Therefore creation must let the receiver earn the good.
 \item This requires free will, struggle, and a world of concealment where effort is
 real.
\end{itemize}

This single concept explains why the world is hard, hidden, and full of choice rather
than a transparent paradise of free gifts.

\subsection{The Desire to Give and the Desire to Receive}

The Ashlag school systematizes the gift structure into two fundamental wills.

\begin{itemize}
 \item \textbf{The desire to bestow}, ratzon lehashpia. This is God's nature, the
 will to give without limit.
 \item \textbf{The desire to receive}, ratzon lekabel. This is the nature of the
 created being, the vessel's will to take in.
\end{itemize}

Creation is, at root, the creation of a desire to receive. The vessel is its wish to
receive the light. The whole spiritual task is to refine that raw receiving so that it
becomes \enquote{receiving in order to bestow,} which restores the vessel to likeness
with the giver. This framing belongs to one influential school, not to every stream of
the tradition.

The contraction, then, is not a single technical step. It is the opening move of a
purpose. God conceals himself so that a free receiver can exist, earn its good, and in
the earning become a giver in turn. The repair that the rest of the chapter describes
is the completion of the purpose that the contraction began.
